\documentclass[11pt]{article}
\input{../../notes-style.tex}

\title{CS 300 -- Chapter 12 Lectures\\\large Sets}
\author{Jose de Lima}
\date{\today}

\begin{document}
\maketitle

\begin{ideabox}
\textbf{The set ADT in one sentence.} A set is a collection of distinct elements with fast membership queries. It's the data structure you reach for any time you care about ``have I seen this before'' or ``how do these two collections compare,'' and almost nothing else. Everything in this chapter is about the operations that make sense on top of that single primitive.
\end{ideabox}

\section{12.1 The Set ADT}

\subsection*{What a set is}

\begin{defnbox}
\textbf{Set.} An unordered collection of distinct elements. No duplicates, no position. $\{3, 7, 9\} = \{9, 7, 3\}$.
\end{defnbox}

The three things a set gives you:
\begin{itemize}[leftmargin=*]
    \item \texttt{add(x)}: insert $x$ if not already present; no-op otherwise.
    \item \texttt{remove(x)}: remove $x$ if present.
    \item \texttt{contains(x)}: membership query.
\end{itemize}

That's it. Everything else (union, intersection, filter, map, subset) is built on top of these three.

\subsection*{Keys vs.\ elements}

When the elements are primitives (ints, strings), each element is its own key. When the elements are objects -- a \texttt{Student} with \texttt{name}, \texttt{phone}, \texttt{id} -- the set uses one field as the key and compares elements by that. The rest of the object is just cargo.

\begin{lstlisting}[basicstyle=\ttfamily\small,frame=single]
// C++ lets you choose: the element type IS the key, or you supply
// a hash/comparator that extracts a key.
std::unordered_set<int> primes;
std::unordered_set<Student, StudentHash, StudentEq> roster;

// Or use a map if you need key/value separation.
std::unordered_map<int, Student> rosterById;
\end{lstlisting}

\begin{notebox}
In Python, sets use \texttt{\_\_hash\_\_} and \texttt{\_\_eq\_\_}; in Java, it's \texttt{hashCode()} and \texttt{equals()}. The contract is always the same: equal elements must hash equal. If you break that contract, the set is not a set -- it's a corrupt data structure waiting to lose an element.
\end{notebox}

\subsection*{Subset testing}

$X \subseteq Y$ iff every element of $X$ is in $Y$. Trivial algorithm:

\begin{lstlisting}[basicstyle=\ttfamily\small,frame=single]
bool isSubset(const std::unordered_set<int>& sub,
              const std::unordered_set<int>& sup) {
    for (int x : sub)
        if (!sup.count(x)) return false;
    return true;
}
\end{lstlisting}

Cost: $O(|X|)$ if \texttt{contains} is $O(1)$ (hash set). If \texttt{sup} is a sorted structure (tree set), it's $O(|X| \log |Y|)$.

\section{12.2 Set Operations}

\subsection*{Union, intersection, difference}

\begin{defnbox}
$X \cup Y$: everything in either set.
$X \cap Y$: everything in both sets.
$X \setminus Y$: everything in $X$ but not in $Y$.
\end{defnbox}

Union and intersection are commutative ($X \cup Y = Y \cup X$). Difference is not.

\begin{examplebox}
$\{54, 19, 75\} \cup \{75, 12\} = \{12, 19, 54, 75\}$ \\
$\{54, 19, 75\} \cap \{75, 12\} = \{75\}$ \\
$\{54, 19, 75\} \setminus \{75, 12\} = \{19, 54\}$ \\
$\{75, 12\} \setminus \{54, 19, 75\} = \{12\}$
\end{examplebox}

\subsection*{C++ implementations}

The \texttt{<algorithm>} header provides set operations on \emph{sorted ranges}, which means any sorted container:

\begin{lstlisting}[basicstyle=\ttfamily\small,frame=single]
#include <set>
#include <algorithm>
#include <iterator>

std::set<int> a{54, 19, 75}, b{75, 12}, out;

std::set_union(a.begin(), a.end(), b.begin(), b.end(),
               std::inserter(out, out.end()));
// out = {12, 19, 54, 75}

out.clear();
std::set_intersection(a.begin(), a.end(), b.begin(), b.end(),
                      std::inserter(out, out.end()));
// out = {75}

out.clear();
std::set_difference(a.begin(), a.end(), b.begin(), b.end(),
                    std::inserter(out, out.end()));
// out = {19, 54}
\end{lstlisting}

For \texttt{std::unordered\_set}, you write the loops explicitly:

\begin{lstlisting}[basicstyle=\ttfamily\small,frame=single]
std::unordered_set<int> setUnion(const std::unordered_set<int>& a,
                                 const std::unordered_set<int>& b) {
    std::unordered_set<int> r(a);
    r.insert(b.begin(), b.end());
    return r;
}

std::unordered_set<int> setIntersect(const std::unordered_set<int>& a,
                                     const std::unordered_set<int>& b) {
    // Iterate the smaller one -- saves time.
    const auto& smaller = a.size() <= b.size() ? a : b;
    const auto& larger  = a.size() <= b.size() ? b : a;
    std::unordered_set<int> r;
    for (int x : smaller)
        if (larger.count(x)) r.insert(x);
    return r;
}

std::unordered_set<int> setDiff(const std::unordered_set<int>& a,
                                const std::unordered_set<int>& b) {
    std::unordered_set<int> r;
    for (int x : a)
        if (!b.count(x)) r.insert(x);
    return r;
}
\end{lstlisting}

\begin{warnbox}
\textbf{Always iterate the smaller set.} For intersection, this shaves the cost from $O(|A|)$ to $O(\min(|A|, |B|))$ -- a huge win when one side is tiny. It's the single most common optimization in real set-heavy code.
\end{warnbox}

\subsection*{Cost analysis}

Assume hash-based sets with $O(1)$ expected insert and lookup.

\begin{center}
\begin{tabular}{lll}
\hline
Operation & Cost & Why \\
\hline
Union $|A| + |B|$ & $O(|A| + |B|)$ & insert all \\
Intersection & $O(\min(|A|, |B|))$ & scan smaller, probe larger \\
Difference $A \setminus B$ & $O(|A|)$ & scan $A$, probe $B$ \\
Subset $A \subseteq B$ & $O(|A|)$ & scan $A$, probe $B$ \\
\hline
\end{tabular}
\end{center}

For sorted (tree) sets, add a $\log n$ factor for each membership probe.

\begin{ideabox}
\textbf{Sorted vs.\ hashed, revisited.} With sorted sets, the two-pointer merge algorithm computes union/intersection/difference in $O(|A| + |B|)$, matching the hashed version without any per-element hash. On small sets that fit in cache, the sorted version can actually be \emph{faster}, because hashing has nontrivial constant overhead. For large sets with large element types, hashed wins. Benchmark before picking.
\end{ideabox}

\subsection*{Filter and map}

Two higher-order operations borrowed from functional programming:

\begin{defnbox}
\textbf{Filter.} Given a predicate $p$, return the subset of elements for which $p$ is true.

\textbf{Map.} Given a function $f$, return the set $\{f(x) : x \in X\}$. The result may be smaller than the input if $f$ is not injective (multiple inputs can map to the same output).
\end{defnbox}

\begin{lstlisting}[basicstyle=\ttfamily\small,frame=single]
#include <unordered_set>
#include <functional>

template<typename T, typename Pred>
std::unordered_set<T> setFilter(const std::unordered_set<T>& s, Pred p) {
    std::unordered_set<T> r;
    for (const auto& x : s) if (p(x)) r.insert(x);
    return r;
}

template<typename T, typename U, typename Fn>
std::unordered_set<U> setMap(const std::unordered_set<T>& s, Fn f) {
    std::unordered_set<U> r;
    for (const auto& x : s) r.insert(f(x));
    return r;
}

// Usage:
auto evens = setFilter(nums, [](int x){ return x % 2 == 0; });
auto lengths = setMap<std::string, size_t>(words,
                   [](const std::string& w){ return w.size(); });
\end{lstlisting}

Python / JavaScript / Rust / Haskell / Scala all have these as first-class methods on their set (and collection) types. C++ got closer in C++20 with \texttt{std::ranges::filter\_view} and \texttt{std::ranges::transform\_view}, but the syntax is still clunky compared to \texttt{xs.filter(...).map(...)} in better-styled languages.

\begin{notebox}
\textbf{Filter is a subset; map is not always.} Filter always returns a subset of the original. Map can shrink the set if collisions occur under $f$: $\{1, -1\} \to \{1\}$ via $f(x) = x^2$. If you need one-to-one semantics, use a list or multiset and track duplicates explicitly.
\end{notebox}

\section{12.3 Static vs.\ Dynamic Sets}

\begin{defnbox}
\textbf{Dynamic set.} Can be modified after construction -- supports add and remove.
\textbf{Static set.} Fixed at construction time. Supports read-only and new-set-producing operations (union, intersection, filter, \ldots), but not in-place add or remove.
\end{defnbox}

\begin{center}
\begin{tabular}{lll}
\hline
Operation & Dynamic & Static \\
\hline
Construct from collection & yes & yes \\
Count elements & yes & yes \\
Search / contains & yes & yes \\
\textbf{Add} & yes & no \\
\textbf{Remove} & yes & no \\
Union / intersection / difference (return new) & yes & yes \\
Filter / map (return new) & yes & yes \\
\hline
\end{tabular}
\end{center}

\subsection*{Why would you ever want a static set?}

\begin{ideabox}
A static set is \emph{immutable}. That buys you several things a mutable set can't:
\begin{itemize}[leftmargin=*]
    \item \textbf{Thread safety for free.} If nothing can change, no locks are needed.
    \item \textbf{Aggressive optimization.} The implementation can use a perfect hash (no collisions), a sorted array (cache-friendly, no per-bucket pointers), or compile-time constants.
    \item \textbf{Safe sharing.} Hand the same set to a dozen consumers and no one can break it.
    \item \textbf{Hashability.} An immutable set can itself be a key in another set or map. (Python's \texttt{frozenset} exists for exactly this.)
\end{itemize}
\end{ideabox}

\subsection*{Choosing in practice}

\begin{center}
\begin{tabular}{ll}
\hline
Use case & Choice \\
\hline
Allowed-country codes loaded at startup & Static -- set once, read forever \\
Stop-words for a search engine & Static -- fixed per language \\
Users currently online & Dynamic -- changes constantly \\
Session IDs issued today & Dynamic -- grows over time \\
Finalized vote tallies (after the election) & Static -- sealed after a deadline \\
Shopping cart contents & Dynamic -- user modifies it \\
\hline
\end{tabular}
\end{center}

\subsection*{C++ idioms for static sets}

C++ doesn't have a built-in frozen set. Common workarounds:

\begin{lstlisting}[basicstyle=\ttfamily\small,frame=single]
// 1. Wrap in const.
const std::unordered_set<std::string> STOP_WORDS{"a", "an", "the", "of"};

// 2. C++17 initializer at static scope.
constexpr std::array<int, 4> ALLOWED_CODES{200, 301, 302, 404};
// Binary-search in a sorted array -- cache-friendly, no dynamic allocation.

// 3. Compile-time perfect hash via gperf, for pathological performance.
\end{lstlisting}

\begin{notebox}
\textbf{Immutability is a spectrum.} ``Const'' in C++ is syntactic -- someone can still \texttt{const\_cast} it away. Python's \texttt{frozenset} is structurally immutable. Rust and Haskell make the distinction first-class. If your codebase has a convention that ``const means truly immutable,'' lean into static sets; if not, you're one \texttt{const\_cast} from heisenbugs.
\end{notebox}

\subsection*{Sets you'll actually use}

\begin{center}
\begin{tabular}{lll}
\hline
Container & Order & Typical implementation \\
\hline
\texttt{std::set} & sorted & red-black tree (ch.\ 9) \\
\texttt{std::unordered\_set} & hash order & open-addressed hash table \\
\texttt{std::multiset} & sorted, allows dups & red-black tree \\
Python \texttt{set} & hash order & open-addressed hash table \\
Python \texttt{frozenset} & hash order & immutable hash table \\
Java \texttt{HashSet} / \texttt{TreeSet} & -- & hash / red-black \\
Rust \texttt{HashSet} / \texttt{BTreeSet} & -- & hash / B-tree \\
\hline
\end{tabular}
\end{center}

\begin{ideabox}
\textbf{Chapter takeaway.} Sets are the deduplicating, order-free, membership-query-oriented sibling of lists. Once you can think of a problem as ``this is a collection where I only care about distinct membership,'' half the algorithm writes itself: iterate one side, probe the other. Pick hash-based when order doesn't matter and you want $O(1)$ membership; pick tree-based when you want sorted iteration or \texttt{lower\_bound}; pick static/frozen when no one should be adding or removing elements. The ADT is simple; the engineering choices it unlocks are not.
\end{ideabox}

\end{document}
